The frozen held-out comparison supports a narrow conclusion: composing the
selected geometry and appearance stages improves source LPIPS and three
target-image endpoints over vanilla V1 under the declared adjusted family.
Surface CD and F-Score remain uncertain despite favorable points. The
selected-versus-Top-1-R comparison is unresolved on every adjusted
endpoint. These results cannot identify semantic matching, matching score,
ray aggregation, or the retargeter as individually necessary. The completed
DEV matched studies also leave their component-level 48-comparison
intervals crossing zero; retrieval coverage and attention allocation are
mechanism evidence rather than generated-quality guarantees.

The outcome-blind panels retain their predetermined samples, while the
conditional Supplement includes joint improvement and separate geometry-
and appearance-loss cases. Those categories expose failures and are not
population success rates. All displayed panels use the scored full-frame
archives and common M1 original-material reference. The native target
conditioning stream remains intact, but the complete generator does not
guarantee pixel-exact target RGB preservation. The PBR selection is a
same-geometry appearance comparison; the evaluated PBR feature/Raw
alternatives use a Top-1 carrier and do not test a shared weighted Top-K
appearance path.

The estimates condition on one frozen full-pipeline checkpoint and one
prescribed generation noise per TEST asset. Two SS training seeds and the
small PBR training diagnostic do not measure full-cascade training-seed
uncertainty. The TEST sources are disjoint within DeformNeXt, but possible
overlap with earlier multi-view pretraining has not been audited. Metrics
use the specified input cameras; they do not establish novel-view editing
quality. FID and KID are renderer-dependent distribution summaries, and
their point differences have no source-cluster interval. Foreground and
GT-crop sensitivity is reported separately without redefining the
full-frame primary endpoints.
